相邻扫描 lane 处理相邻候选，子空间优先存储使其编码访问相邻；LUT 分解则独立减少每查询的查找表。

\subsection{地址映射与四字节打包}

令 $o_\ell$ 为倒排列表 $\ell$ 的起点，$r$ 为列表内行号，$p=o_\ell+r$ 为物理位置。若 $Z[m,p]$ 表示主码子空间 $m$ 的八比特编码，逐字节扫描副本将其存于字节偏移 $mN+p$，而候选优先布局使用 $pM+m$。当各 lane 访问同一列表内连续位置时，前者使固定 $m$ 的访问相邻。跨列表边界时，位置可能跳跃，因此不保证每个 warp 都对应一次连续事务。

打包表示将同一候选的四个子空间编码组成一个无符号字：
\begin{equation}
 \mathsf{w}[j,p]=\sum_{t=0}^{3}2^{8t}Z[4j+t,p],\qquad
 \operatorname{offset}_{\rm bytes}(\mathsf{w}[j,p])=4(jN+p).
 \label{eq:5-packed-layout}
\end{equation}
通过 $(\mathsf{w}[j,p]\gg(8t))\mathbin{\&}255$ 精确恢复四个编码。每个编码再提供\cref{eq:5-lut-factorisation}中的高、低半字节索引。打包四个子空间，与拆分一个子空间的八比特标签不同，也不等于打包四个候选的编码。

\begin{figure}[tbp]
\centering
% Address schematic; four lanes stand for a subset of a warp.
\begin{tikzpicture}[font=\small,
 box/.style={draw=ThesisNavy!65,rounded corners=2pt,align=center,
   text width=2.6cm,minimum height=1.2cm,inner sep=4pt},
 arr/.style={-{Latex[length=2mm]},draw=ThesisNavy!80,thick}]
\foreach \a in {0,1,2,3} {
 \node[box,fill=ThesisNavy!5] (w\a) at (3.3*\a,0)
  {lane $\a$\\$\mathsf{w}[j,p+\a]$};
}
\node[align=center] at (4.95,1.05)
 {Fixed $j$: adjacent 32-bit words across candidate positions};
\node[box,text width=11.5cm,minimum height=1.6cm] (decode) at (4.95,-2.5)
 {Unpack each lane's word independently at its position $p_a$\\
  $Z[4j,p_a]\;|\;Z[4j+1,p_a]\;|\;Z[4j+2,p_a]\;|\;Z[4j+3,p_a]$\\
  Four code bytes $\longrightarrow$ shifts and masks $\longrightarrow$ LUT indices};
\draw[arr] (w0.south) -- (decode.north west);
\draw[arr] (w1.south) -- (decode.north);
\draw[arr] (w2.south) -- (decode.north);
\draw[arr] (w3.south) -- (decode.north east);
\end{tikzpicture}

\caption[打包后的候选并行访问]{打包后的候选并行访问，图中仅画出部分 warp lane。固定字组 $j$ 时，连续候选位置访问相邻四字节 word。每个 word 包含同一候选的四个子空间编码，解包保持编码值，随后再进行 LUT 拆分。此图说明地址关系，不表示实测内存事务数量。}
\label{fig:5-packed-access}
\end{figure}

建库路径明确将 $[M,N]$ 字节数组重新打包为 $[M/4,N]$ word 数组。仅把原字节数组重新解释为 \texttt{uint32\_t}，会合并相邻候选的字节，而不是需要的子空间。实现要求 $M$ 能被四整除，并拒绝不完整的最后一个 word。重排在建库时读写一次 $NM$ 编码字节，旧存储和新存储短暂共存；它位于搜索计时边界之外。

\subsection{打包访问省下什么、保持什么}

一次 word load 取得四个编码，随后在寄存器中移位与掩码解包。
32 个连续位置的逻辑窗口为 128 字节，实际事务取决于对齐、活跃 lane 及内存层次。
载荷仍为每候选 $M$ 字节。两种被测布局已按子空间组织，因此开关对照测试的是
打包／加载／解包的综合效果，没有单独分离转置、指令或 cache 的贡献。

\Cref{sec:7-3}比较打包路径与逐字节路径的吞吐量、召回率及波动。
